% ============================================================
%  CHAPTER 9 — Agent Orchestration and Skill Library
% ============================================================
\chapter{Agent Orchestration and Skill Library}
\label{ch:agent}

\section{Overview}

The \texttt{EmbodiedAgent} class (\texttt{embodied\_agent\_v2.py}) is the top-level
orchestrator that integrates all subsystems into a coherent perceive--reason--act
loop.
It manages the user interface (voice and text), coordinates with the VLM reasoner and
perception manager, executes the selected skill, and enforces the safety policy.

\section{Initialisation Sequence}

The agent initialises all sub-components sequentially on startup, with timeout-guarded
connection to each ROS action server:

\begin{enumerate}
  \item \textbf{VLMReasoner}: validate Google API key, connect to Gemini API.
  \item \textbf{PerceptionManager}: subscribe to camera topics, connect to YOLO service.
  \item \textbf{StateManager}: subscribe to \texttt{/joint\_states} and
        \texttt{/mobile\_base\_controller/odom}.
  \item \textbf{Load skills}: instantiate all skill objects and connect their
        action clients.
  \item \textbf{TTS}: connect to \texttt{/tts} action server (10 s timeout).
  \item \textbf{STT}: enumerate microphones and select best device.
  \item \textbf{Publishers}: advertise \texttt{/agent/camera\_annotated} and
        \texttt{/agent/status}.
  \item \textbf{Safety handlers}: register SIGINT, SIGTERM, and ROS shutdown hooks.
\end{enumerate}

\section{Skill Library}

\subsection{Abstract Base Class}

All skills inherit from \texttt{BaseSkill} and implement three mandatory methods:

\begin{lstlisting}[language=Python, caption={BaseSkill abstract interface.}]
class BaseSkill:
    def check_affordance(self, params, state) -> (bool, str):
        """Check if skill can execute given current state.
        Returns (can_execute: bool, reason: str)."""

    def execute(self, params) -> bool:
        """Execute the skill. Return True on success."""

    def get_description(self) -> str:
        """Return description string shown to VLM."""
\end{lstlisting}

The separation between \texttt{check\_affordance} and \texttt{execute} is critical:
it allows the agent to verify preconditions \emph{before} attempting physical
execution, and to communicate failures back to the VLM without having moved the robot.

\subsection{Skill: grab\_bottle}

\begin{description}[noitemsep]
  \item[Description shown to VLM:] \textit{grab\_bottle() — Grasp any detected
    object (bottle, phone, cup, mug, etc.) in front of the robot using vision and
    arm control.}
  \item[Affordance check:]
    (1) Gripper must be \texttt{empty}.
    (2) At least one graspable object must appear in detected objects.
    (3) The reach script must exist on disk.
  \item[Execution:] Calls \texttt{reach\_object\_v5\_torso\_descent\_working.py}
    as a subprocess with a 3-minute timeout.
    Returns success if the subprocess exits with code 0.
\end{description}

\subsection{Skill: search\_with\_head}

\begin{description}[noitemsep]
  \item[Description shown to VLM:] \textit{search\_with\_head(object\_name=`bottle')
    — Move head to scan surroundings for object.}
  \item[Affordance check:] Always available.
  \item[Execution:] Sweeps head through 12 pan/tilt positions (Table~\ref{tab:head_positions}
    in Chapter~\ref{ch:manipulation}). Runs YOLO detection at each position after
    a 2 s settle delay. Returns \texttt{True} when target is found.
\end{description}

\subsection{Skill: handover\_to\_person}

\begin{description}[noitemsep]
  \item[Description shown to VLM:] \textit{handover\_to\_person(object=`bottle')
    — Detect person with YOLO+depth, extend arm toward them, and hand over held
    object. Use after grab\_bottle when the goal is to deliver an object to someone.}
  \item[Affordance check:]
    (1) Gripper must be \texttt{holding:*} (non-empty).
    (2) A \texttt{person} must appear in detected objects.
  \item[Execution:] Nine-step pipeline — home, clear planning scene, get TF,
    detect person, add person collision box, move arm to handover height,
    speak, open gripper, go home.
\end{description}

\subsection{Skill: push\_aside\_object}

\begin{description}[noitemsep]
  \item[Description shown to VLM:] \textit{push\_aside\_object(blocking\_object=\textquoteleft cup\textquoteright)
    --- Push a blocking object sideways to clear the path to the target.
    Use when an object obstructs the arm's path to the grasp target.}
  \item[Affordance check:]
    (1) A detected object matching the blocking name must be visible.
    (2) A 3D position must be available (depth + TF).
  \item[Execution:] Determines lateral push direction (away from robot centreline),
    positions the closed gripper alongside the object, then executes a Cartesian
    lateral push via MoveIt.
    Respects a maximum table-lateral limit (0.55\,m) to prevent pushing objects
    off the edge.
\end{description}

\subsection{Skill: navigate\_to}

\begin{description}[noitemsep]
  \item[Description shown to VLM:] \textit{navigate\_to(location=\textquoteleft table\textquoteright)
    --- Drive the robot base to a named location.}
  \item[Affordance check:] The named location must exist in \texttt{config/locations.yaml}.
  \item[Execution:] Calls \texttt{NavigationClient.navigate\_to(location\_name)}
    which sends a goal to the \texttt{move\_base} action server.
    Named locations (home, table, shelf) are defined in \texttt{config/locations.yaml}
    as (x, y, $\theta$) waypoints in the map frame.
\end{description}

\subsection{Simple Motion Skills}

Four skills wrap PAL \texttt{play\_motion} pre-defined trajectories:

\begin{table}[H]
\centering
\caption{Simple motion skills.}
\label{tab:simple_skills}
\begin{tabular}{lll}
\toprule
\textbf{Skill} & \textbf{Motion name} & \textbf{Affordance} \\
\midrule
\texttt{go\_home}   & \texttt{home}  & Always \\
\texttt{wave}       & \texttt{wave}  & Always \\
\texttt{open\_hand} & \texttt{open}  & Always \\
\texttt{close\_hand}& \texttt{close} & Always (warns if already holding) \\
\bottomrule
\end{tabular}
\end{table}

All four share a single \texttt{play\_motion} action client instance (class-level
singleton) to avoid creating multiple connections to the same server.

\section{Multi-Step Execution Loop}

The \texttt{execute\_command()} method implements the iterative control loop.
Figure~\ref{fig:exec_loop} shows the state machine.

\begin{figure}[H]
\centering
\begin{tikzpicture}[
  s/.style={rectangle, rounded corners=4pt, draw=darkblue, thick,
            fill=lightblue!30, minimum width=3.8cm, minimum height=0.7cm,
            align=center, font=\small},
  d/.style={diamond, draw=darkblue, thick, fill=yellow!20,
            minimum width=2.2cm, minimum height=0.8cm,
            align=center, font=\footnotesize},
  a/.style={-{Stealth[length=5pt]}, thick, draw=darkblue},
  node distance=0.5cm and 1.8cm
]
  \node[s] (cap)    {Capture image\\+ run YOLO};
  \node[s, below=of cap] (state)  {Assemble state};
  \node[s, below=of state](query) {Query VLM};
  \node[d, below=of query](done)  {skill\\=done?};
  \node[d, below=of done] (aff)   {Affordance\\OK?};
  \node[s, below=of aff]  (exec)  {Execute skill};
  \node[d, below=of exec] (suc)   {Success?};
  \node[s, right=2.5cm of suc]   (fail)  {Ensure safe\\pose};
  \node[s, right=2.5cm of done]  (finish){Task\\complete};

  \draw[a] (cap) -- (state);
  \draw[a] (state) -- (query);
  \draw[a] (query) -- (done);
  \draw[a] (done.east) -- node[above, font=\tiny]{yes} (finish.west);
  \draw[a] (done.south) -- node[right, font=\tiny]{no} (aff.north);
  \draw[a] (aff.west)  -- ++(-1.2,0) |- node[left, font=\tiny]{no — retry with failure ctx} (query.west);
  \draw[a] (aff.south) -- node[right, font=\tiny]{yes} (exec.north);
  \draw[a] (exec) -- (suc);
  \draw[a] (suc.east)  -- node[above, font=\tiny]{no}  (fail.west);
  \draw[a] (suc.west)  -- ++(-1.5,0) |- node[left, font=\tiny]{yes — next step} (cap.west);
\end{tikzpicture}
\caption{Multi-step execution loop state machine.}
\label{fig:exec_loop}
\end{figure}

Key design choices in the loop:

\begin{itemize}
  \item \textbf{Affordance failure as context}: when a skill's affordance check fails,
        the failure reason is passed back to the VLM in the next query, allowing it to
        select a different skill (\eg search\_with\_head if grab fails because object
        is not visible).

  \item \textbf{Loop detection}: if the same skill is selected twice consecutively,
        the loop assumes convergence and stops.

  \item \textbf{Visualisation}: at each step, the annotated image and status JSON
        are published to \texttt{/agent/camera\_annotated} and \texttt{/agent/status}
        for the web dashboard.
\end{itemize}

\section{State Manager}

The \texttt{StateManager} maintains a live representation of the robot and scene state,
updated from ROS subscribers.

\subsection{State Dictionary}

At each step, the following state dict is assembled and passed to the VLM:

\begin{lstlisting}[language=JSON, caption={State dictionary passed to VLM at each step.}]
{
  "gripper":           "holding:bottle",
  "base_location":     "table",
  "detected_objects":  ["bottle at (185, 290)", "person at (410, 200)"],
  "task_history":      ["\u2713 search_with_head", "\u2713 grab_bottle"],
  "spatial_relations": ["bottle left_of person", "bottle near person"]
}
\end{lstlisting}

\subsection{Scene Graph and Symbolic Rule Engine}

Spatial relations are computed from 3D object positions in the \texttt{base\_footprint}
frame (X\,=\,forward, Y\,=\,left, Z\,=\,up), extracted via depth camera back-projection
and TF transforms.
These geometric base facts seed the \textbf{SymbolicRuleEngine} (see
Section~\ref{sec:rule_engine}), which derives higher-order facts by forward chaining
to fixpoint.

\textbf{Geometric base facts (3D):}
\begin{itemize}[noitemsep]
  \item \textbf{left\_of/right\_of}: $|A_Y - B_Y| > 0.08$\,m (hysteresis band).
  \item \textbf{above/below}: $|A_Z - B_Z| > 0.08$\,m.
  \item \textbf{near\_3d}: Euclidean distance $< 0.30$\,m.
  \item \textbf{reachable}:  < 1.20$\,m and $|A_Y| < 0.65$\,m (arm envelope).
  \item \textbf{on\_table}:  < 0.90$\,m and  < 1.5$\,m.
\end{itemize}

Self-referential pairs ( = B$) are filtered to prevent nonsensical facts such as
\texttt{left\_of(bottle,\,bottle)}.
If depth data is unavailable, the system falls back to 2D pixel heuristics with
a \texttt{relation\_source: 2d} flag logged for diagnostics.

\textbf{Derived facts (from rule engine):}
\begin{itemize}[noitemsep]
  \item \texttt{can\_grasp(A)}: \texttt{reachable(A) $\wedge$ graspable(A)}
  \item \texttt{handover\_possible(X,P)}: \texttt{holding(gripper,X) $\wedge$ near\_3d(robot,P)}
  \item \texttt{co\_located(A,B)}: objects on same surface
  \item \texttt{stacked\_on(A,B)}: \texttt{above(A,B) $\wedge$ near\_xy(A,B)}
\end{itemize}

\textbf{VLM semantic enrichment:}
After each VLM call, the \texttt{semantic\_facts} field of the JSON response is
merged into the fact base and the rule engine re-fires.
This allows the VLM to assert visual observations that geometry cannot compute
(e.g., \texttt{bottle\_is\_open}, \texttt{person\_is\_reaching}), completing the
neuro-symbolic loop.

\subsection{Gripper State Tracking}

Gripper status is tracked at two levels:
\begin{enumerate}[noitemsep]
  \item \textbf{Automatic}: the \texttt{/joint\_states} subscriber monitors
        \texttt{gripper\_left\_finger\_joint} position. Position $<$ 5 mm indicates
        open; $>$ 15 mm indicates closed.
  \item \textbf{Explicit}: skills call \texttt{state.update\_gripper\_status()} after
        a successful grasp or release, setting \texttt{"holding:bottle"} or
        \texttt{"empty"} respectively.
\end{enumerate}

Explicit updates take precedence and are not overridden by the automatic tracker,
preventing false positives when the gripper closes on air.

\section{Speech Interface}

\subsection{Speech-to-Text}

The agent supports two STT backends with automatic fallback:

\begin{enumerate}[noitemsep]
  \item \textbf{Whisper (local)}: \texttt{recognizer.recognize\_whisper(audio,
        model="base")}. Preferred: runs offline, handles accented speech well,
        $\approx$74M parameters, runs in real time on CPU.
  \item \textbf{Google STT}: \texttt{recognizer.recognize\_google(audio)}. Fallback
        when Whisper fails or is unavailable.
\end{enumerate}

The microphone listens for up to 8 seconds with a 15-second phrase time limit.
Ambient noise calibration is performed for 0.3 seconds before each listen.

\subsection{Text-to-Speech}

The agent speaks status messages at key points in the loop via the PAL TTS server.
Messages are designed to be informative and short:

\begin{center}
\begin{tabular}{ll}
\toprule
\textbf{Event} & \textbf{Spoken message} \\
\midrule
Start of step & ``Looking at the scene.'' \\
Objects found & ``I can see: bottle, person.'' \\
VLM thinking  & ``Thinking about what to do.'' \\
VLM rationale & First 200 chars of rationale \\
Skill start   & ``Executing grab bottle.'' \\
Skill done    & ``Grab bottle done.'' \\
Task complete & ``Task complete!'' \\
Failure       & ``Grab bottle failed. Returning to safe position.'' \\
Safe mode     & ``Too many failures. Entering safe mode.'' \\
\bottomrule
\end{tabular}
\end{center}

\section{Safety System}

\subsection{Affordance Gating}

Every skill is gated by \texttt{check\_affordance()} before execution.
This prevents physically impossible or dangerous actions — for example, attempting
to grasp when the gripper is already holding an object, or attempting a handover
when no person is detected.

\subsection{Consecutive Failure Tracking}

The agent tracks the number of consecutive skill failures:

\begin{lstlisting}[language=Python, caption={Failure counter and safe mode trigger.}]
if not success:
    self.consecutive_failures += 1
    if self.consecutive_failures >= self.max_consecutive_failures:  # = 3
        self._enter_safe_mode(
            reason="{} consecutive failures".format(self.consecutive_failures))
        return {'success': False, 'safe_mode': True, ...}
    self._ensure_safe_pose(reason="skill_failure")
\end{lstlisting}

A successful skill resets the counter to zero.

\subsection{Safe Mode}

When 3 consecutive failures are detected, the agent enters \textbf{safe mode}:

\begin{enumerate}[noitemsep]
  \item \texttt{open\_hand} is executed (release any held object).
  \item \texttt{go\_home} is executed (return arm to rest).
  \item All further command execution is blocked.
  \item The user must send the command \texttt{reset} to resume.
\end{enumerate}

\subsection{Shutdown Sequence}

On SIGINT, SIGTERM, or ROS shutdown, a registered handler executes:
\begin{enumerate}[noitemsep]
  \item \texttt{go\_home} — return arm to rest.
  \item \texttt{open\_hand} — release any held object.
\end{enumerate}

This ensures the robot is in a stable, safe configuration regardless of what it was
doing when interrupted.

\section{Live Visualisation}

At each control step, the agent publishes:

\begin{itemize}[noitemsep]
  \item \texttt{/agent/camera\_annotated} (\texttt{sensor\_msgs/Image}): the current
        RGB frame annotated with YOLO bounding boxes and the active skill name overlaid
        as a text banner.
  \item \texttt{/agent/status} (\texttt{std\_msgs/String}): a JSON string with the
        current skill, gripper state, detected objects, location, and last 5 tasks.
\end{itemize}

These are consumed by \texttt{web\_visualiser.py}, providing a real-time browser
dashboard for monitoring the agent during operation.
